\section{Randomized --- Zero P}
%%% generated by the blueprint pipeline from TCSlib.Complexity.Randomized.ZeroP
\subsection*{Randomized --- Zero P: Introduction}

Declarations of the Lean module \texttt{TCSlib.Complexity.Randomized.ZeroP}.

\subsection*{Randomized --- Zero P: Results}

\begin{definition}[Pair-encoded list step]
\lean{Randomized.PIT.ptStep}
\label{Randomized.PIT.ptStep}
\leanok
\uses{Complexity.pairFstD, Complexity.pairSndD, Randomized.dropPrefixByLen, Randomized.takePrefixByLen, Turing.pairEncode}
For a Boolean list $d$, define a Boolean list by applying the first-component projection
to $d$ as the first entry of a pair encoding. The second entry is another pair encoding:
its first component is the suffix obtained by dropping a prefix of length specified by
the encoding of the first-component projections of the two projections of $d$; its
second component is the second projection of the second projection of $d$, followed by
an encoding of the corresponding prefix.
\end{definition}

\begin{theorem}[Iteration of the point-testing step]
\lean{Randomized.PIT.iterate_ptStep}
\label{Randomized.PIT.iterate_ptStep}
\leanok
\uses{Complexity.BinArith.encWords, Complexity.BinArith.encWords_append_singleton, Complexity.pairFstD_pairEncode, Complexity.pairSndD_pairEncode, Randomized.PIT.ptStep, Randomized.dropPrefixByLen, Randomized.takePrefixByLen, Turing.pairEncode}
\difficulty{4}
For every block length $b\in\mathbb{N}$, bit string $u$, list $P$ of bit strings, and
iteration count $i\in\mathbb{N}$, iterating the point-testing step $i$ times from the
encoded state with counter consisting of $b$ ones, remaining string $u$, and stored list
$P$ yields the state with the same counter, remaining string obtained by deleting the
first $ib$ bits of $u$, and stored list formed by appending, for each $j<i$, the
length-$b$ block beginning at position $jb$ of $u$.
\end{theorem}

\begin{theorem}[Polynomial-time computability of the loop]
\lean{Randomized.PIT.polyTimeComputable_ptLoop}
\label{Randomized.PIT.polyTimeComputable_ptLoop}
\leanok
\uses{Complexity.BinArith.forLoop, Complexity.BinArith.forLoop_bound_of_ctx, Complexity.BinArith.polyTimeComputable_forLoop, Complexity.BinArith.polyTimeComputable_slot2, Complexity.BinArith.polyTimeComputable_slot3, Complexity.PolyTimeComputable, Complexity.PolyTimeComputable.append, Complexity.PolyTimeComputable.comp, Complexity.PolyTimeComputable.pairEncode, Complexity.length_pair_components_le, Complexity.pairFstD, Complexity.pairFstD_pairEncode, Complexity.pairSndD, Complexity.pairSndD_pairEncode, Complexity.polyTimeComputable_const, Complexity.polyTimeComputable_pairFstD, Randomized.PIT.ptStep, Randomized.dropPrefixByLen, Randomized.polyTimeComputable_dropPrefixByLen, Randomized.polyTimeComputable_takePrefixByLen, Randomized.takePrefixByLen, Turing.length_pairEncode, Turing.pairEncode}
\difficulty{5}
The function on bit strings obtained by iterating the defined loop step is computable in
polynomial time: there exist a finite-state Turing machine and natural numbers $C,c$
such that, on every input $x$, the machine halts after exactly $C(|x|+1)^c$ steps and
outputs the result of applying the loop to $x$.
\end{theorem}

\begin{definition}[First projected component]
\lean{Randomized.PIT.qL}
\label{Randomized.PIT.qL}
\leanok
\uses{Complexity.pairFstD}
For a Boolean list $Q$, define its associated list to be the first component obtained by
applying the given pair projection to $Q$.
\end{definition}

\begin{definition}[Nested projected component]
\lean{Randomized.PIT.qN}
\label{Randomized.PIT.qN}
\leanok
\uses{Complexity.pairFstD, Complexity.pairSndD}
For a Boolean list $Q$, define its associated list to be the first component obtained by
applying the pair projection to the second component of $Q$.
\end{definition}

\begin{definition}[Third nested projected component]
\lean{Randomized.PIT.qG₁}
\leanok
\uses{Complexity.pairFstD, Complexity.pairSndD}
For a Boolean list $Q$, define its associated list to be the first component obtained by
applying the pair projection to the second component of the second component of $Q$.
\end{definition}

\begin{definition}[Fourth nested projected component]
\lean{Randomized.PIT.qG₂}
\leanok
\uses{Complexity.pairFstD, Complexity.pairSndD}
For a Boolean list $Q$, define its associated list to be the first component obtained by
applying the pair projection to the second component of the second component of the
second component of $Q$.
\end{definition}

\begin{definition}[Fifth nested second component]
\lean{Randomized.PIT.qP}
\label{Randomized.PIT.qP}
\leanok
\uses{Complexity.pairSndD}
For a Boolean list $Q$, let $q_P(Q)$ be the list obtained by taking the second component
in the pairing scheme five times in succession.
\end{definition}

\begin{definition}[Modular evaluation trial test]
\lean{Randomized.PIT.trialTest}
\statementsource{Vad12}{
  pdf-bf2ce5422088-p014-b001,
  pdf-bf2ce5422088-p016-b003
}
\label{Randomized.PIT.trialTest}
\leanok
\uses{Complexity.BinArith.beq, Randomized.PIT.evalModW, Randomized.PIT.qG₁, Randomized.PIT.qG₂, Randomized.PIT.qN, Randomized.PIT.qP, Turing.pairEncode}
For Boolean lists $Q$ and $R$, the trial test returns the Boolean result of comparing,
via the binary arithmetic equality test, the modular evaluations of two encoded
expressions. Each evaluation uses an encoding built from the extracted length data of
$Q$, the list $R$, the fifth nested second component of $Q$, and, respectively, the
first or second extracted generator from $Q$.
\end{definition}

\begin{definition}[Trial state transition]
\lean{Randomized.PIT.trStep}
\label{Randomized.PIT.trStep}
\leanok
\uses{Complexity.pairFstD, Complexity.pairSndD, Randomized.PIT.qL, Randomized.PIT.trialTest, Randomized.dropPrefixByLen, Randomized.takePrefixByLen, Turing.pairEncode}
For a Boolean list $d$, form a new encoded state whose first component is the first
component extracted from $d$. Its second component encodes the list obtained by dropping
from a length encoding the prefix specified by the extracted length data and first
component of the second component of $d$, together with a singleton Boolean list. That
Boolean is the conjunction of the first component of the third nested second component
of $d$ and the trial test applied to the first component of $d$ and the corresponding
prefix of the length encoding.
\end{definition}

\begin{theorem}[Iteration of the trial-testing step]
\lean{Randomized.PIT.iterate_trStep}
\label{Randomized.PIT.iterate_trStep}
\leanok
\uses{Complexity.pairFstD_pairEncode, Complexity.pairSndD_pairEncode, Randomized.PIT.qL, Randomized.PIT.trStep, Randomized.PIT.trialTest, Randomized.dropPrefixByLen, Randomized.takePrefixByLen, Turing.pairEncode}
\difficulty{4}
For bit strings $Q,s$, a bit $b$, and an iteration count $i\in\mathbb{N}$, iterating the
trial-testing step $i$ times from the encoded state with query $Q$, remaining string
$s$, and flag $b$ yields the state with query $Q$, remaining string obtained by deleting
the first $i|q|$ bits of $s$, and flag equal to $b$ conjoined with the outcomes of the
trial test on each of the first $i$ consecutive blocks of $s$ of length $|q|$, where $q$
is the query's length component.
\end{theorem}

\begin{theorem}[Polynomial-Time Computability of the Trial Test]
\lean{Randomized.PIT.polyTimeComputable_trialTest}
\label{Randomized.PIT.polyTimeComputable_trialTest}
\leanok
\uses{Complexity.BinArith.beq, Complexity.BinArith.polyTimeComputable_app, Complexity.BinArith.polyTimeComputable_beq, Complexity.BinArith.polyTimeComputable_slot2, Complexity.BinArith.polyTimeComputable_slot3, Complexity.BinArith.polyTimeComputable_slot4, Complexity.BinArith.polyTimeComputable_test_app, Complexity.PolyTimeComputable, Complexity.PolyTimeComputable.comp, Complexity.PolyTimeComputable.pairEncode, Complexity.pairFstD, Complexity.pairSndD, Complexity.polyTimeComputable_pairFstD, Complexity.polyTimeComputable_pairSndD, Randomized.PIT.evalModW, Randomized.PIT.polyTimeComputable_evalModW, Randomized.PIT.qG₁, Randomized.PIT.qG₂, Randomized.PIT.qN, Randomized.PIT.qP, Randomized.PIT.trialTest, Turing.pairEncode}
\difficulty{5}
The Boolean function that takes a bit string $z$, applies the total first and second
projections of the pair encoding to obtain $Q$ and $R$, and returns the one-bit string
containing the trial-test result on $(Q,R)$ is computable by a deterministic Turing
machine in time bounded by a polynomial in the input length.
\end{theorem}

\begin{theorem}[Polynomial-Time Computability of the Transition Loop]
\lean{Randomized.PIT.polyTimeComputable_trLoop}
\label{Randomized.PIT.polyTimeComputable_trLoop}
\leanok
\uses{Complexity.BinArith.forLoop, Complexity.BinArith.forLoop_bound_of_ctx, Complexity.BinArith.polyTimeComputable_app, Complexity.BinArith.polyTimeComputable_forLoop, Complexity.BinArith.polyTimeComputable_slot2, Complexity.BinArith.polyTimeComputable_slot3, Complexity.BinArith.polyTimeComputable_test_app, Complexity.PolyTimeComputable, Complexity.PolyTimeComputable.pairEncode, Complexity.length_pair_components_le, Complexity.pairFstD, Complexity.pairFstD_pairEncode, Complexity.pairSndD, Complexity.pairSndD_pairEncode, Complexity.polyTimeComputable_and, Complexity.polyTimeComputable_headD, Complexity.polyTimeComputable_pairFstD, Randomized.PIT.polyTimeComputable_trialTest, Randomized.PIT.qL, Randomized.PIT.trStep, Randomized.PIT.trialTest, Randomized.dropPrefixByLen, Randomized.polyTimeComputable_dropPrefixByLen, Randomized.polyTimeComputable_takePrefixByLen, Randomized.takePrefixByLen, Turing.length_pairEncode, Turing.pairEncode}
\difficulty{5}
The function that applies the transition step repeatedly to an input bit string for its
length plus one iterations is polynomial-time computable.
\end{theorem}

\begin{definition}[First extracted component]
\lean{Randomized.PIT.pN}
\label{Randomized.PIT.pN}
\leanok
\uses{Complexity.pairFstD}
For a Boolean list $z$, let $p_N(z)$ be its first component in the pairing scheme.
\end{definition}

\begin{definition}[Second extracted component]
\lean{Randomized.PIT.pI}
\label{Randomized.PIT.pI}
\leanok
\uses{Complexity.pairFstD, Complexity.pairSndD}
For a Boolean list $z$, let $p_I(z)$ be the first component of the second component of
$z$ in the pairing scheme.
\end{definition}

\begin{definition}[Third extracted component]
\lean{Randomized.PIT.pG₁}
\leanok
\uses{Complexity.pairFstD, Complexity.pairSndD}
For a Boolean list $z$, let $p_{G_1}(z)$ be the first component of the third nested
second component of $z$ in the pairing scheme.
\end{definition}

\begin{definition}[Fourth extracted component]
\lean{Randomized.PIT.pG₂}
\leanok
\uses{Complexity.pairFstD, Complexity.pairSndD}
For a Boolean list $z$, let $p_{G_2}(z)$ be the first component of the fourth nested
second component of $z$ in the pairing scheme.
\end{definition}

\begin{definition}[Fifth extracted component]
\lean{Randomized.PIT.pR}
\label{Randomized.PIT.pR}
\leanok
\uses{Complexity.pairSndD}
For a Boolean list $z$, let $p_R(z)$ be the fifth nested second component of $z$ in the
pairing scheme.
\end{definition}

\begin{definition}[Padded Boolean list]
\lean{Randomized.PIT.pM}
\label{Randomized.PIT.pM}
\leanok
\uses{Randomized.PIT.pI}
For a Boolean list $z$, define its padded form by appending twelve copies of `true` to
the list $p_I(z)$.
\end{definition}

\begin{definition}[Point sequence encoding]
\lean{Randomized.PIT.pPts}
\label{Randomized.PIT.pPts}
\leanok
\uses{Complexity.BinArith.forLoop, Randomized.PIT.pI, Randomized.PIT.pN, Randomized.PIT.pR, Randomized.PIT.ptStep, Turing.pairEncode}
For a Boolean list $z$, define its point sequence encoding by applying the iteration
`forLoop` with step `ptStep` to the nested pair encoding of $p_N(z)$, $p_I(z)$ with four
`true` entries appended, and $p_R(z)$ paired with the empty list.
\end{definition}

\begin{definition}[Boolean pipeline output]
\lean{Randomized.PIT.pipe}
\label{Randomized.PIT.pipe}
\leanok
\uses{Complexity.BinArith.forLoop, Complexity.pairFstD, Complexity.pairSndD, Randomized.PIT.pG₁, Randomized.PIT.pG₂, Randomized.PIT.pM, Randomized.PIT.pN, Randomized.PIT.pPts, Randomized.PIT.trStep, Turing.pairEncode}
For a Boolean list $z$, define the pipeline output to be the first entry of the list
obtained by applying `forLoop` with step `trStep` to the encoded data built from eight
concatenated copies of $p_M(z)$, the nested encodings of two copies of $p_M(z)$,
$p_N(z)$, $p_{G_1}(z)$, $p_{G_2}(z)$, and the specified projections of $pPts(z)$. The
resulting list is projected three times by `pairSndD` before taking its first entry,
with `false` as the default if it is empty.
\end{definition}

\begin{theorem}[Polynomial-time computability of the pipeline output]
\lean{Randomized.PIT.polyTimeComputable_pipe}
\label{Randomized.PIT.polyTimeComputable_pipe}
\leanok
\uses{Complexity.BinArith.forLoop, Complexity.BinArith.polyTimeComputable_app, Complexity.BinArith.polyTimeComputable_slot2, Complexity.BinArith.polyTimeComputable_slot3, Complexity.BinArith.polyTimeComputable_slot4, Complexity.BinArith.polyTimeComputable_test_app, Complexity.PolyTimeComputable, Complexity.PolyTimeComputable.append, Complexity.PolyTimeComputable.comp, Complexity.PolyTimeComputable.pairEncode, Complexity.pairFstD, Complexity.pairSndD, Complexity.polyTimeComputable_const, Complexity.polyTimeComputable_headD, Complexity.polyTimeComputable_pairFstD, Complexity.polyTimeComputable_pairSndD, Randomized.PIT.pG₁, Randomized.PIT.pG₂, Randomized.PIT.pI, Randomized.PIT.pM, Randomized.PIT.pN, Randomized.PIT.pPts, Randomized.PIT.pR, Randomized.PIT.pipe, Randomized.PIT.polyTimeComputable_ptLoop, Randomized.PIT.polyTimeComputable_trLoop, Randomized.PIT.ptStep, Randomized.PIT.trStep, Turing.pairEncode}
\difficulty{6}
The Boolean function computed by the specified pipeline on finite bit strings has a
polynomial-time algorithm: there are a deterministic multi-tape Turing machine and
natural numbers $C,c$ such that, on every input $z$, the machine halts within
$C(|z|+1)^c$ steps and outputs the one-bit string containing the pipeline’s Boolean
value on $z$.
\end{theorem}

\begin{definition}[Consecutive Boolean blocks]
\lean{Randomized.PIT.pts}
\label{Randomized.PIT.pts}
\leanok
For natural numbers $n,m$ and a Boolean list $r$, form a list of $n$ Boolean lists by
taking, for each $i=0,\ldots,n-1$, the block of length $m+4$ beginning at position
$i(m+4)$ in $r$.
\end{definition}

\begin{definition}[Indexed Boolean list block]
\lean{Randomized.PIT.rawj}
\label{Randomized.PIT.rawj}
\leanok
For natural numbers $n,m,j$ and a Boolean list $r$, define the $j$th raw block to be the
list obtained by first removing the initial $n(m+4)$ entries of $r$, then removing the
next $j\,2(m+12)$ entries, and taking up to $2(m+12)$ entries from what remains.
\end{definition}

\begin{theorem}[Pipeline test characterization]
\lean{Randomized.PIT.pipe_eq}
\label{Randomized.PIT.pipe_eq}
\leanok
\uses{Complexity.BinArith.beq, Complexity.BinArith.encWords, Complexity.BinArith.forLoop_pairEncode, Complexity.LogProg.bitsVal, Complexity.pairFstD_pairEncode, Complexity.pairSndD_pairEncode, Randomized.PIT.ctxOf, Randomized.PIT.evalModW, Randomized.PIT.iterate_ptStep, Randomized.PIT.iterate_trStep, Randomized.PIT.pG₁, Randomized.PIT.pG₂, Randomized.PIT.pI, Randomized.PIT.pM, Randomized.PIT.pN, Randomized.PIT.pPts, Randomized.PIT.pR, Randomized.PIT.pipe, Randomized.PIT.pts, Randomized.PIT.qG₁, Randomized.PIT.qG₂, Randomized.PIT.qL, Randomized.PIT.qN, Randomized.PIT.qP, Randomized.PIT.rawj, Randomized.PIT.trialTest, Turing.pairEncode}
\difficulty{6}
For all $n,m\in\mathbb{N}$ and bit strings $G_1,G_2,r$, the pipeline applied to the
encoding of $n$, $m$, $G_1$, $G_2$, and $r$ returns true exactly when, for every integer
$j$ with $0\le j<8(m+12)$, the two gate evaluations modulo the modulus specified by the
$j$th raw block of $r$ agree. In each evaluation, the context is formed from $n$, that
raw block, and the list of input points specified by the corresponding blocks of $r$,
and the gate lists are $G_1$ and $G_2$.
\end{theorem}

\begin{theorem}[Taking after dropping from a truncated list]
\lean{Randomized.PIT.take_drop_take}
\label{Randomized.PIT.take_drop_take}
\leanok
\difficulty{3}
Let $l$ be a finite bit string and let $N,a,b$ be natural numbers with $a+b\le N$.
Taking $b$ entries after dropping the first $a$ entries from the first $N$ entries of
$l$ gives the same list as taking $b$ entries after dropping the first $a$ entries of
$l$.
\end{theorem}

\begin{theorem}[Zero test characterization]
\lean{Randomized.PIT.zeroTest_eq}
\label{Randomized.PIT.zeroTest_eq}
\leanok
\uses{Complexity.LogProg.bitsVal, Randomized.AlgCircuit, Randomized.AlgCircuit.ZeroTestCore, Randomized.AlgCircuit.eval, Randomized.AlgCircuit.size, Randomized.AlgCircuit.testModulus, Randomized.AlgCircuit.testPoint, Randomized.AlgCircuit.zeroTest, Randomized.PIT.ptVal, Randomized.PIT.pts, Randomized.PIT.rawj, Randomized.PIT.take_drop_take, Randomized.blockVal, Randomized.inputBits, Randomized.modParam, Randomized.numTrials, Randomized.testLen}
\difficulty{5}
Let $D$ be an algebraic circuit on $n$ inputs, and let $r$ be a finite bit string. Its
zero-test on $r$ is true exactly when, for every integer $j$ with $0 \le j <
8(\operatorname{size}(D)+12)$, evaluating $D$ at the point whose $i$th coordinate is one
plus the integer encoded by the $(\operatorname{size}(D)+4)$-bit block of $r$ beginning
at position $i(\operatorname{size}(D)+4)$, with arithmetic in the residue ring modulo
one plus the integer encoded by the $2(\operatorname{size}(D)+12)$-bit block of $r$
beginning after its first $n(\operatorname{size}(D)+4)$ bits and then at offset
$j\,2(\operatorname{size}(D)+12)$, gives zero.
\end{theorem}

\begin{theorem}[Zero Value of the Encoded Evaluation]
\lean{Randomized.PIT.bitsVal_evalModW_nil}
\label{Randomized.PIT.bitsVal_evalModW_nil}
\leanok
\uses{Complexity.LogProg.bitsVal, Randomized.PIT.ctxOf, Randomized.PIT.evalModW, Randomized.PIT.gateRun_eq, Randomized.PIT.runVal_encState, Turing.pairEncode}
\difficulty{3}
For every $n\in\mathbb{N}$, bit string $R$, and list $P$ of bit strings, evaluating the
encoded context formed from $n$, $R$, and $P$ with an empty second component yields a
bit string whose binary value is zero.
\end{theorem}

\begin{theorem}[Vanishing of a natural residue]
\lean{Randomized.PIT.natCast_eq_zero_iff_of_lt}
\label{Randomized.PIT.natCast_eq_zero_iff_of_lt}
\leanok
\difficulty{4}
For natural numbers $v$ and $k$ with $v<k$, the residue of $v$ in $\mathbb Z/k\mathbb Z$
is zero if and only if $v=0$.
\end{theorem}

\begin{definition}[Zero-test predicate]
\lean{Randomized.PIT.zeropW}
\statementsource{Vad12}{
  pdf-bf2ce5422088-p016-b002,
  pdf-bf2ce5422088-p016-b003
}
\label{Randomized.PIT.zeropW}
\leanok
\uses{Complexity.BinArith.uHead, Complexity.BinArith.uRest, Complexity.pairFstD, Complexity.pairSndD, Randomized.PIT.pipe, Randomized.PIT.sI, Randomized.PIT.scanRun, Randomized.PIT.validW, Turing.pairEncode}
For a Boolean list $z$, define the zero-test predicate to hold exactly when either
`validW` fails on the first projection of $z$, or `pipe` returns false on the nested
encoding whose entries are the first projection's unary head, the result of applying
`sI` to `scanRun` of that projection, its unary remainder, the empty list, and the
second projection of $z$.
\end{definition}

\begin{theorem}[Polynomial-Time Computability of the Zero Predicate]
\lean{Randomized.PIT.polyTimeComputable_zeropW}
\label{Randomized.PIT.polyTimeComputable_zeropW}
\leanok
\uses{Complexity.BinArith.polyTimeComputable_app, Complexity.BinArith.polyTimeComputable_slot3, Complexity.BinArith.polyTimeComputable_test_app, Complexity.BinArith.polyTimeComputable_uHead, Complexity.BinArith.polyTimeComputable_uRest, Complexity.BinArith.uHead, Complexity.BinArith.uRest, Complexity.PolyTimeComputable, Complexity.PolyTimeComputable.comp, Complexity.PolyTimeComputable.pairEncode, Complexity.pairFstD, Complexity.pairSndD, Complexity.polyTimeComputable_const, Complexity.polyTimeComputable_not, Complexity.polyTimeComputable_or, Complexity.polyTimeComputable_pairFstD, Complexity.polyTimeComputable_pairSndD, Randomized.PIT.polyTimeComputable_pipe, Randomized.PIT.polyTimeComputable_scanRun, Randomized.PIT.polyTimeComputable_validW, Randomized.PIT.sI, Randomized.PIT.scanRun, Randomized.PIT.validW, Randomized.PIT.zeropW, Turing.pairEncode}
\difficulty{5}
The Boolean function that sends each bit string $z$ to the one-bit string containing the
zero-predicate result on $z$ is computable by a deterministic Turing machine in time
bounded by a polynomial in the input length.
\end{theorem}

\begin{theorem}[Agreement of Two Zero-Test Verifiers]
\lean{Randomized.PIT.zeropW_pairEncode}
\label{Randomized.PIT.zeropW_pairEncode}
\leanok
\uses{Complexity.BinArith.uHead, Complexity.BinArith.uHead_encodeNat, Complexity.BinArith.uRest, Complexity.BinArith.uRest_encodeNat, Complexity.pairFstD_pairEncode, Complexity.pairSndD_pairEncode, Randomized.AlgCircuit, Randomized.AlgCircuit.WellFormed, Randomized.AlgCircuit.decode, Randomized.AlgCircuit.encode, Randomized.AlgCircuit.encodeGates, Randomized.AlgCircuit.encode_of_decode, Randomized.AlgCircuit.size, Randomized.AlgCircuit.zeroTest, Randomized.PIT.bitsVal_evalModW_nil, Randomized.PIT.evalModW_spec, Randomized.PIT.natCast_eq_zero_iff_of_lt, Randomized.PIT.pipe_eq, Randomized.PIT.pts, Randomized.PIT.rawj, Randomized.PIT.sI_scanRun, Randomized.PIT.validW, Randomized.PIT.validW_iff, Randomized.PIT.zeroTest_eq, Randomized.PIT.zeropW, Randomized.zeropVerifier, Turing.pairEncode}
\difficulty{6}
For all finite bit strings $x$ and $r$, the word-level verifier on the paired encoding
of $x$ and $r$ returns the same Boolean value as the circuit-level verifier on $x$ and
random string $r$.
\end{theorem}

\begin{theorem}[Efficiency of the polynomial zero test verifier]
\lean{Randomized.eff_zeropVerifier}
\label{Randomized.eff_zeropVerifier}
\leanok
\uses{Randomized.PIT.polyTimeComputable_zeropW, Randomized.PIT.zeropW_pairEncode, Randomized.boolVerifier, Randomized.eff_boolVerifier_of_test, Randomized.polyTimeModel, Randomized.zeropVerifier}
\difficulty{1}
The verifier that parses an algebraic circuit, accepts if parsing fails or the circuit
is not well formed, and otherwise rejects exactly when the circuit’s randomized zero
test accepts runs in polynomial time in the verifier model.
\end{theorem}

\begin{theorem}[Zero-polynomial testing in polynomial-time coRP]
\lean{Randomized.inCoRP_ZEROP_polyTime}
\statementsource{Vad12}{
  pdf-bf2ce5422088-p022-b004,
  pdf-bf2ce5422088-p017-b001,
  pdf-bf2ce5422088-p017-b002
}
\label{Randomized.inCoRP_ZEROP_polyTime}
\leanok
\uses{Randomized.InCoRP, Randomized.ZEROP, Randomized.eff_zeropVerifier, Randomized.inCoRP_ZEROP, Randomized.polyTimeModel}
\difficulty{1}
The language of encodings of nonempty, well-formed arithmetic circuits over the integers
that compute the zero polynomial belongs to coRP for the polynomial-time verifier model.
Equivalently, there is a polynomial-time randomized verifier that rejects every such
zero circuit with certainty and accepts every string outside this language with
probability at least $2/3$.
\end{theorem}

\begin{theorem}[Polynomial identity testing for zero is in BPP]
\lean{Randomized.inBPP_ZEROP_polyTime}
\label{Randomized.inBPP_ZEROP_polyTime}
\leanok
\uses{Randomized.InBPP, Randomized.InCoRP.inBPP_polyTime, Randomized.ZEROP, Randomized.inCoRP_ZEROP_polyTime, Randomized.polyTimeModel}
\difficulty{1}
The language of encodings of well-formed arithmetic circuits over the integers whose
associated multivariate polynomial is zero belongs to BPP.
\end{theorem}

\begin{theorem}[Zero-polynomial circuits lie in coRP]
\lean{Randomized.ZEROP_mem_coRP_PTM}
\statementsource{Vad12}{
  pdf-bf2ce5422088-p022-b004,
  pdf-bf2ce5422088-p017-b001,
  pdf-bf2ce5422088-p017-b002
}
\label{Randomized.ZEROP_mem_coRP_PTM}
\leanok
\uses{Randomized.ZEROP, Randomized.coRP_PTM, Randomized.inCoRP_ZEROP_polyTime, Randomized.mem_coRP_PTM_iff}
\difficulty{1}
The language of encodings of well-formed arithmetic circuits over the integers whose
computed polynomial is zero belongs to coRP.
\end{theorem}

\begin{theorem}[Zero Polynomial Testing in Bounded Error Polynomial Time]
\lean{Randomized.ZEROP_mem_BPP_PTM}
\statementsource{Vad12}{
  pdf-bf2ce5422088-p016-b003,
  pdf-bf2ce5422088-p017-b001,
  pdf-bf2ce5422088-p017-b002,
  pdf-bf2ce5422088-p023-b001
}
\label{Randomized.ZEROP_mem_BPP_PTM}
\leanok
\uses{Randomized.BPP_PTM, Randomized.ZEROP, Randomized.inBPP_ZEROP_polyTime, Randomized.mem_BPP_PTM_iff}
\difficulty{1}
The language of encodings of well-formed algebraic circuits over the integer polynomials
that compute the zero polynomial has a randomized polynomial-time decision procedure: it
halts within a polynomial bound on every choice sequence and outputs the correct
membership bit with probability at least $2/3$ on every input.
\end{theorem}
